is this element of @@M0@@ independent of the particular choice of the @@M1@@ with intersection @@M2@@?

\sgaStableTarget{sga6.e14.sec.7}\section*{7. Chern classes of perfect complexes}

\sgaStableTarget{sga6.e14.para.7-1}\textbf{7.1.} In seeking to develop a cohomological Riemann--Roch formula, with or without denominators, beyond the rather artificial hypothesis that the schemes under consideration admit an ample invertible Module, one is necessarily confronted with the question of defining, for a perfect complex of Modules @@M3@@ on a scheme @@M4@@, Chern classes
@@M5@@
generalizing the Chern classes of locally free Modules in SGA 5 VII. The same question also arises for general locally ringed topoi (compare [10]) and for perfect complexes on paracompact spaces ringed by continuous complex functions, working then with ordinary integral cohomology. For a compact space, the @@M6@@ constructed with perfect complexes is known to be the same as the ``naive'' one constructed with locally free Modules (Exp. IV), and in this case one is trivially reduced to the classical definition. But even for a locally compact space the problem posed seems nontrivial, since the integral cohomology of such a space is not always the inverse limit of the cohomologies of its compact subsets.

\sgaStableTarget{sga6.e14.para.7-2}\textbf{7.2.} More precisely, for a closed subset @@M7@@ of @@M8@@ such that @@M9@@ is zero (in the sense of derived categories), i.e. such that the cohomology sheaves of @@M10@@ have supports in @@M11@@, we would like to define Chern classes \uline{with support in} @@M12@@
@@M13@@
which would recover the ``ordinary'' (sic) Chern classes of 7.1 by means of the canonical homomorphism
@@M14@@
